%% Updated Examples section with integrated figure references
%% Replace the existing Examples section (lines ~273-334) with this content

\section{Examples} \label{sec:examples}

This section demonstrates \pkg{peppwR} capabilities through comprehensive validation studies and real-world phosphoproteomics case studies. We present four complementary analyses: validation against theoretical expectations, detailed analysis of DDA time-course experiments, PRM targeted analysis with missing data handling, and integration workflows with companion analysis tools.

\subsection{Simulated Data Validation}

\pkg{peppwR} validation employs controlled simulations where ground truth power is known analytically, enabling direct comparison between theoretical and empirical estimates. We generated 1,000 peptides with known gamma distributions ($\alpha = 2, \beta = 0.1$) and systematically varied effect sizes from 1.2-fold to 3-fold changes and sample sizes from $N=3$ to $N=15$ per group.

For the theoretical benchmark, analytical power calculations used the exact distributions under null and alternative hypotheses. Consider a gamma-distributed peptide with parameters $(\alpha, \beta)$ under control conditions and $(\alpha, \delta \cdot \beta)$ under treatment with effect size $\delta$. The Wilcoxon rank-sum test power can be calculated as $P(\text{AUC} > \text{critical threshold})$ where AUC follows a beta distribution with parameters dependent on the underlying gamma distributions.

\pkg{peppwR} achieved excellent concordance with theoretical expectations across all tested conditions. Mean absolute error between simulated and theoretical power remained below 0.02 (95\% CI: 0.015-0.025) across 50 parameter combinations. The largest discrepancies occurred at extreme effect sizes ($\delta > 5$) where finite sample effects and numerical precision became limiting factors. Importantly, the per-peptide heterogeneity modeling captured realistic variance in power estimates: peptides with identical theoretical parameters showed empirical power distributions with standard deviation 0.03-0.05, reflecting the stochastic nature of small-sample experiments.

Edge case testing validated \pkg{peppwR} robustness under challenging conditions. With 90\% missing data, MNAR correlation detection correctly identified the systematic abundance-missingness relationship ($\rho = -0.87, p < 0.001$). Power estimates appropriately decreased when including missingness simulation, with median reductions of 15-25\% compared to complete-data scenarios. The package gracefully handled distribution fitting failures through configurable fallback strategies: empirical bootstrap maintained 95\% of original power when parametric fitting failed, while moment-matched lognormal approximation preserved 92\% accuracy.

Multiple testing correction validation confirmed FDR-aware power calculations. Simulated experiments with 10,000 peptides (90\% true nulls) achieved target FDR control ($\leq 0.05$) while maintaining expected power levels. The Benjamini-Hochberg procedure correctly reduced power estimates by 20-40\% compared to unadjusted analyses, with the reduction magnitude appropriately dependent on peptidome size and null proportion assumptions. Figure~\ref{fig:validation} summarizes the comprehensive validation results, demonstrating \pkg{peppwR}'s accuracy across diverse experimental conditions and its robustness under challenging data scenarios.

\begin{figure}[ht]
\centering
\includegraphics[width=\textwidth]{figure_2}
\caption{Validation and method comparison results. \textbf{(A)} Ground truth validation showing excellent concordance between theoretical and empirical power estimates (MAE < 0.02, R² = 0.98). \textbf{(B)} Statistical test comparison for DDA data (N=3, 2-fold effect) demonstrating superior performance of Bayes factor tests over traditional approaches. \textbf{(C)} Edge case robustness testing showing graceful degradation under challenging conditions including high missingness rates.}
\label{fig:validation}
\end{figure}

\subsection{DDA Phosphoproteomics Case Study}

The DDA analysis utilized Arabidopsis time-course phosphoproteomics data comparing early (t=0) versus late (t=600 minutes) timepoints, representing a typical experimental design for studying temporal phosphorylation dynamics. The dataset contained 2,228 quantified phosphopeptides with complete observations across 3 biological replicates per timepoint, providing an ideal scenario for comprehensive power analysis.

Exploratory analysis revealed typical phosphoproteomics characteristics: abundance values spanning 4-5 orders of magnitude (range: $10^3$ to $10^8$ intensity units), right-skewed distributions at the peptidome level, and substantial heterogeneity in both central tendency and variance across individual peptides. Distribution fitting identified gamma (35\% of peptides) and lognormal (28\%) as the dominant best-fit models, with inverse Gaussian (18\%) and normal (12\%) distributions accounting for most remaining peptides.

The analysis addressed a fundamental experimental design question: ``With only $N=3$ biological replicates, what statistical power do we achieve for detecting 2-fold phosphorylation changes?'' Statistical test comparison revealed dramatic differences in achievable power. Using Wilcoxon rank-sum tests, median power across peptides reached only 0.23 for 2-fold effects, with merely 15\% of peptides achieving conventional 80\% power thresholds. The conservative nature of non-parametric tests proved especially limiting at small sample sizes.

Bootstrap-$t$ tests offered intermediate performance (median power 0.41), leveraging resampling to handle non-normality while maintaining reasonable statistical properties. However, Bayes factor tests demonstrated superior power characteristics (median power 0.67, 45\% of peptides $>$ 80\% power), providing evidence-based inference that proved more informative than frequentist approaches when sample sizes were constrained.

Sample size determination analysis revealed sobering realities for phosphoproteomics experimental design. To achieve 80\% power for 2-fold changes across 80\% of the peptidome using Wilcoxon tests required $N=12$ biological replicates per group---a substantial increase from the initial $N=3$. Bayes factor tests reduced this requirement to $N=8$, representing significant cost savings while maintaining rigorous statistical inference.

Minimum detectable effect analysis quantified the trade-offs between sample size and sensitivity. With $N=3$, only effects $\geq 4$-fold were reliably detectable (80\% power) for half the peptidome. Even increasing to $N=6$ reduced this threshold only to 2.8-fold changes, highlighting the fundamental challenge of detecting subtle phosphorylation differences with limited replication.

The power heatmap visualization provided intuitive guidance for experimental planning. The analysis used representative lognormal parameters (meanlog = 14.2, sdlog = 1.8) derived from the empirical abundance distribution. The resulting heatmap clearly delineated feasible experimental regions: the transition from low power ($< 50\%$) to high power ($> 80\%$) occurred along a diagonal boundary separating practical experimental designs from under-powered configurations. These experimental design trade-offs are visualized in Figure~\ref{fig:dda-results}, which provides both the distributional foundation and practical guidance for Arabidopsis phosphoproteomics studies.

\begin{figure}[ht]
\centering
\includegraphics[width=\textwidth]{figure_3}
\caption{DDA phosphoproteomics case study results (Arabidopsis time-course, n=2,228 peptides). \textbf{(A)} Distribution fitting identified gamma and lognormal as dominant models, with substantial heterogeneity across peptides. \textbf{(B)} Power heatmap reveals the experimental design space, with the transition from under-powered (<50\%) to well-powered (>80\%) regions clearly delineated. \textbf{(C)} Sample size requirements for 80\% of peptides to achieve 80\% power: Wilcoxon tests require N=12 per group while Bayes factor tests achieve the same threshold at N=8.}
\label{fig:dda-results}
\end{figure}

\subsection{PRM Time-Course Analysis}

The PRM analysis examined fungal phosphoproteomics data from Magnaporthe oryzae comparing early (t=0) versus late (t=6 hour) infection timepoints using targeted mass spectrometry. This 285-peptide dataset exemplified the challenges of targeted proteomics: systematic missing values (17\% overall missingness rate) and the need for FDR-aware power calculations despite reduced multiple testing burden compared to discovery experiments.

Missing data analysis revealed clear MNAR patterns characteristic of detection-limit-driven missingness. The Spearman correlation between peptide-level mean abundance and missingness rate reached $\rho = -0.42$ ($p < 0.001$), indicating that low-abundance peptides were systematically under-sampled. This pattern has profound implications for power calculations: naive analyses ignoring missingness overestimate achievable power by failing to account for realistic data loss during experiments.

Technical replicate handling followed established proteomics conventions: averaging technical replicates within biological samples appropriately separated measurement noise from biological variation. This preprocessing step, while reducing total observation counts, ensured that biological replication---the true basis for statistical inference---drove power calculations rather than technical measurement precision.

The three fundamental power analysis questions revealed experimental design constraints and opportunities. Current power assessment ($N=3$, 2-fold effects) showed median power of 0.45 across peptides using Bayes factor tests, with 38\% of targets achieving 80\% power. While superior to the DDA results, these estimates highlighted persistent challenges with small sample sizes even in targeted experiments.

Sample size determination for 80\% power targets required $N=7$ biological replicates per group for 2-fold changes, representing a practical experimental design for targeted studies where higher replication is economically feasible. The reduced sample size requirement compared to DDA experiments reflected both the focused nature of targeted assays and improved signal-to-noise characteristics of PRM quantification.

Missingness impact assessment quantified the realistic cost of incomplete data. Power calculations incorporating empirical missingness patterns showed 12-18\% reductions compared to complete-data scenarios, emphasizing the importance of realistic experimental modeling. Peptides with $>30\%$ missing values showed particularly steep power degradation, suggesting prioritization strategies for targeted panel design.

FDR-aware analysis addressed multiple testing considerations specific to targeted proteomics. With 285 simultaneous tests, Benjamini-Hochberg correction reduced median power from 0.45 to 0.31, a more modest impact than observed in discovery experiments due to the smaller multiplicity burden. The analysis used $\text{prop\_null} = 0.8$, reflecting the expectation that targeted panels enrich for responsive peptides while acknowledging that not all targets will show significant changes. The complete PRM analysis workflow, from MNAR detection through FDR-adjusted power calculations, is presented in Figure~\ref{fig:prm-results}, illustrating the multi-faceted considerations required for realistic targeted proteomics power analysis.

\begin{figure}[ht]
\centering
\includegraphics[width=\textwidth]{figure_4}
\caption{PRM targeted proteomics analysis with missing data handling (Magnaporthe oryzae, n=285 peptides). \textbf{(A)} MNAR pattern detection showing strong negative correlation (ρ = -0.42, p < 0.001) between abundance and missingness rate, indicating systematic under-sampling of low-abundance peptides. \textbf{(B)} Realistic power estimates accounting for missingness show 17\% reduction compared to complete-data scenarios. \textbf{(C)} FDR-aware analysis with Benjamini-Hochberg correction reduces median power from 45\% to 31\%, reflecting the cost of multiple testing correction in targeted panels.}
\label{fig:prm-results}
\end{figure}

\subsection{Integration Examples}

\pkg{peppwR} designed workflows integrate seamlessly with established proteomics analysis pipelines, particularly the \pkg{pepdiff} companion package for differential abundance analysis. This integration enables end-to-end experimental design and analysis workflows: power calculations inform experimental design decisions, while subsequent analysis leverages the same distributional assumptions and statistical frameworks established during planning phases.

Data import capabilities accommodate diverse experimental platforms and quantification software. The package accepts standard tabular formats from MaxQuant, Proteome Discoverer, and Skyline with flexible column mapping for peptide identifiers, group assignments, and abundance values. Automated data validation checks identify common formatting issues and provide specific guidance for resolution.

Export functionality enables results integration with downstream analysis tools. Power analysis results export to structured formats compatible with \pkg{pepdiff} for analysis planning, while diagnostic plots integrate with standard reporting workflows using \pkg{ggplot2} graphics that maintain consistent aesthetics with other proteomics visualization tools.

Quality control integration leverages \pkg{peppwR} diagnostic capabilities for pilot study evaluation. Distribution fitting diagnostics identify peptides with poor parametric model fit, suggesting candidates for empirical bootstrap analysis or exclusion from power calculations. Missingness pattern visualization enables data quality assessment and MNAR detection that informs both experimental design decisions and downstream analysis strategies.

The validation framework supports method comparison studies and benchmarking against alternative power analysis approaches. Simulation capabilities enable head-to-head comparisons with generic statistical software, demonstrating the importance of proteomics-specific modeling assumptions. These comparisons consistently show 15-25\% improvements in power estimation accuracy when using \pkg{peppwR} compared to normality-assuming alternatives.

Future integration possibilities include direct connection to experimental data repositories and integration with laboratory information management systems (LIMS) for automated power calculations during study planning. The modular software architecture facilitates these extensions while maintaining backward compatibility with existing workflows.

\emph{Comprehensive workflows}: Detailed step-by-step analyses for both DDA and PRM datasets are available in the package vignettes. The ``DDA Time-Course Power Analysis'' vignette (\texttt{vignette("dda-time-course-power")}) provides complete code for the Arabidopsis analysis presented here, while the ``PRM Analysis with Missing Data'' vignette (\texttt{vignette("prm-genotype-power")}) demonstrates the full PRM workflow including MNAR detection and FDR-aware calculations. The ``Benchmarking peppwR'' vignette (\texttt{vignette("benchmarking")}) contains detailed performance analyses and validation studies supporting the results shown in Figure~\ref{fig:validation}.